\documentclass[10pt,a4paper]{amsart}
\usepackage[margin=19mm]{geometry}
\usepackage{amsmath,amssymb,mathtools}
\usepackage[scr=rsfs,cal=euler]{mathalfa}
\usepackage{xfrac}
\usepackage[unicode,hidelinks,bookmarks=false]{hyperref}
\newcommand{\ie}{i.e.\ }
\newcommand{\cf}{cf.\ }
\newcommand{\resp}{respectively\ }
\newcommand{\Subsection}{\subsection}
\providecommand{\nobreakdash}{\nobreak\hbox{-}}
\setlength{\emergencystretch}{2em}
\multlinegap0pt
\usepackage{bm}
\usepackage{bbm}
\usepackage{dsfont}
\usepackage{xspace}
\usepackage{cancel}
\usepackage{mathtools}
\usepackage[shortlabels]{enumitem}
\usepackage{amsthm}
\makeatletter
\@ifundefined{thm}{\newtheorem{thm}{Theorem}}{}
\@ifundefined{lem}{\newtheorem{lem}[thm]{Lemma}}{}
\makeatother
\newcommand\N{\mathbb N}
\newcommand\ass {\operatorname{Ass}}
\newcommand\fp{\mathfrak p}
\newcommand\fq{\mathfrak q}
\newcommand\supp {\operatorname{Supp}}
\title[A new characterization of the weakly Laskerian (FSF) modules]{A new characterization of the weakly Laskerian (FSF) modules}
\author{Ali Fathi}
\date{Source: arXiv:2507.04573v1 (CC BY 4.0)}
\begin{document}
\maketitle

\noindent\emph{Source and licence.} A new characterization of the weakly Laskerian (FSF) modules. Version arXiv:2507.04573v1 (CC BY 4.0), distributed under the Creative Commons Attribution 4.0 International licence, as recorded in the arXiv licence metadata for this exact version and shown on the abstract page https://arxiv.org/abs/2507.04573v1. Full rights notice: licenses/arxiv-2507.04573-CC-BY-4.0.txt.\par\smallskip

\noindent\textit{Editorial scope.} Counted unit: the Thm whose source label is \texttt{None} in the article (source0.tex lines 114--120, source label \texttt{None}), delivered together with the article's own complete original proof of it (source0.tex lines 121--152, one \texttt{proof} environment of 32 lines). No mathematics is written, repaired or re-derived: statement and proof are the article's own text line for line. Required context retained uncounted: none. Every definition, display and label of those lines is retained unchanged. Omitted from this package and not redistributed: 1--113, 153--185. Nothing inside a retained excerpt is omitted or altered: every retained body is delivered exactly as the article's lines are written, so no omission receipt of any kind accompanies this package. Citations: the retained excerpts carry no citation command at all, so no bibliography entry of the article is transcribed and no reference is redistributed. Reference adaptation: none was needed and none was made; every reference inside the delivered unit resolves inside it, so no unresolved reference or citation is produced. Printed slips kept literal: no printed word, symbol, hypothesis or number of the counted statement or of its proof was altered. Layout and class: the article's own class is \texttt{amsart} with packages including amsfonts, amsmath, amstext, amssymb, amsthm, mathrsfs, amscd, enumerate, verbatim, xy; the wrapper is the lane's pinned \texttt{amsart} editorial wrapper at 10pt on A4, which restates the theorem-like environments the retained text needs and adds the article's own macro definitions that the retained text needs (\texttt{\textbackslash N}, \texttt{\textbackslash ass}, \texttt{\textbackslash fp}, \texttt{\textbackslash fq}, \texttt{\textbackslash supp}), with extra wrapper packages (bm, bbm, dsfont, xspace, cancel, mathtools, enumitem). The delivered numbering is the unit's own. Assets: no retained span contains a figure environment, a graphics-inclusion command, a PDF-page inclusion, a table or a TikZ picture; no image file is copied into this package, the package declares no asset dependency and no image file is redistributed with it. Licence: version arXiv:2507.04573v1 of this article is distributed under the Creative Commons Attribution 4.0 International licence, as recorded in the arXiv licence metadata for this exact version and shown on the abstract page https://arxiv.org/abs/2507.04573v1. A full rights notice is delivered at licenses/arxiv-2507.04573-CC-BY-4.0.txt. Characters: every retained excerpt is pure ASCII, so the delivered engine, XeLaTeX with its default Latin Modern text font, reports no missing character.
\begin{thm}Let $M$ be an $R$-module. Then the following statements are equivalent:
\begin{enumerate}[\rm(i)]
\item $M$ is a weakly Laskerian $R$-module;
\item there is a finitely generated submodule $N$ of $M$ such that $\supp_R(M/N)=\ass_R(M/N)$ and $\supp_R(M/N)$ is a finite set;
\item $M$ is an FSF $R$-module.
\end{enumerate}
\end{thm}

\begin{proof}
The implication (ii)$\Rightarrow$(iii) is clear. To prove the implication (iii)$\Rightarrow$(i), assume that $N$ is a finitely generated submodule  of $M$ such that $\supp_R(M/N)$ is finite. Suppose that $L$ is an arbitrary submodule of $M$. It follows from the exact sequence
$$0\rightarrow(N+L)/L\rightarrow M/L\rightarrow M/(N+L)\rightarrow 0$$
that $\ass_R(M/L)\subseteq\ass_R((M/N)/((N+L)/N))\cup\ass_R(N/(N\cap L))$. Therefore  $\ass_R(M/L)$ is a finite set and consequently $M$ is weakly Laskerian.

(i)$\Rightarrow$(ii). Assume that $M$ is weakly Laskerian and to prove (ii) it is sufficient for us to show that  $\supp_R(M/N)=\ass_R(M/N)$ for some finitely generated submodule $N$ of $M$.

We set  $M_0:=0$ (the zero submodule of $M$). Assume $i\in\N$.  If $\supp_R(M/M_{i-1})=\ass_R(M/M_{i-1})$, then we set $N:=M_{i-1}$ and we end the process.  Otherwise,  we construct $x_i, M_i,  \fq_i, \fp_i, \Sigma_i$ as follows. We set $\Sigma_i:=\supp_R(M/M_{i-1})\setminus\ass_R(M/M_{i-1})$. Since $\Sigma_i$ is a non-empty set of ideals of $R$ and $R$ is Noetherian,  $\Sigma_i$ has a maximal element under inclusion, say  $\fp_i$.
  There exist $\fq_i\in\ass_R(M/M_{i-1})$ and $x_i\in M$ such that $\fq_i=(0:_Rx_i+M_{i-1})\subset\fp_i$ (note that $\fq_i\in\ass_R(M/M_{i-1})$ while $\fp_i\notin\ass_R(M/M_{i-1})$ and so $\fq_i\neq\fp_i$). Now, we  set $M_i:=M_{i-1}+\fp_ix_i$. It is clear that $M_i$ is a finitely generated submodule of $M$. We claim that after a finite number of steps the procedure must stop.  Assume for the sake of contradiction that $\supp_R(M/M_i)\neq\ass_R(M/M_i)$ for all $i$ and $x_i, M_i,  \fq_i, \fp_i, \Sigma_i$ are constructed for all $i\in\N$ as above. Before continuing the proof, we need the following two lemmas.

  \begin{lem}\label{lem1} For each $ j\in \N$, $\ass_R(M_j/M_{j-1})=\{\fq_j\}$.
  \end{lem}
  \begin{proof}
  Since $\fp_j\nsubseteq\fq_j$, we obtain $M_{j-1}\subset M_j$ and so $\ass_R(M_j/M_{j-1})\neq\emptyset$. Now, suppose that  $\fp\in\ass_R(M_j/M_{j-1})$. Then $\fp=(0:_Ra_jx_j+M_{j-1})$ for some $a_j\in\fp_j$. Therefore
   $\fp a_j\subseteq\fq_j$ and so $\fp\subseteq \fq_j$. The reverse inclusion is clear and hence  $\fp=\fq_j$.
  \end{proof}

  \begin{lem} If  $1\leq i\leq j$, then $\fp_i=(0:_Rx_i+M_j)\in\ass_R(M/M_j)$ and $\fp_i\nsubseteq\fp_j$ whenever $1\leq i<j$.
  \end{lem}
  \begin{proof}
  First, we prove that $\fp_j=(0:_Rx_j+M_j)$ for all $j\in\N$. Assume that $r\in(0:_Rx_j+M_j)$. Therefore $rx_j=m_{j-1}+a_jx_j$ for some $m_{j-1}\in M_{j-1},a_j\in\fp_j$. It follows that $r-a_j\in\fq_j$ and so $r\in\fp_j$. This shows that $(0:_Rx_j+M_j)\subseteq\fp_j$. The reverse inclusion is clear and so $\fp_j=(0:_Rx_j+M_j)\in\ass_R(M/M_j)$ for all $1\leq j$.

  Next, we show that    if $1\leq i<j$, then $\fp_i\nsubseteq\fp_j$. Assume for the sake of contradiction that $\fp_i\subseteq\fp_j$ for some $i, j$ with $i<j$. We can assume that $j$ is the least positive integer with the property  that there exists $i$ with $i<j$ such that $\fp_i\subseteq\fp_j$. If $\fp_i=\fp_j$, then since $\fp_i=(0:_Rx_i+M_i)$ we have $\fp_j\in\ass_R(M/M_i)$ while $\fp_j\notin\ass_R(M/M_{j-1})$ by definition. Therefore $i<j-1$ and  there exists $i+1\leq k< j$ such that $\fp_j\in\ass_R(M/M_{k-1})$ and $\fp_j\notin\ass_R(M/M_{k})$. Also, if $\fp_i\subset \fp_j$,  then, by the maximality of $\fp_i$ in $\Sigma_i$,  $\fp_j\in\ass_R(M/M_{i-1})$ while  $\fp_j \notin\ass_R(M/M_{j-1})$ by definition. Therefore there exits $i\leq k <j$ such that $\fp_j\in\ass_R(M/M_{k-1})$ and $\fp_j\notin\ass_R(M/M_k)$. In both cases,
  we deduce from the exact sequence
 $$0\rightarrow M_k/M_{k-1}\rightarrow M/M_{k-1}\rightarrow M/M_k\rightarrow 0$$
 that $\fp_j\in\ass_R(M_k/M_{k-1})$ and so Lemma \ref{lem1} implies that $\fp_j=\fq_k\subset\fp_k$. Since $\fp_i\subseteq\fp_j$, we obtain $\fp_i\subset\fp_k$, which is impossible by the minimality of $j$. Therefore $\fp_i\nsubseteq\fp_j$ for all $1\leq i<j$.

 Finally, we prove by induction on $j$  that $\fp_i=(0:_Rx_i+M_j)$ for all $1\leq i\leq j$. If $j=1$, then we have $\fp_1=(0:_Rx_1+M_1)$. Now assume that $j>1$ and the result has been proved for the smaller values of $j$.  The case $i=j$ is proved at the beginning of the proof of this lemma. So assume  that $i<j$. By the above proof, we have $\fp_i\nsubseteq \fp_j$ and hence there exists $s\in \fp_i\setminus\fp_j$. Suppose that $r\in(0:_Rx_i+M_j)$. Therefore $rx_i=m_{j-1}+a_jx_j$ for some $m_{j-1}\in M_{j-1}, a_j\in\fp_j$. Hence
$srx_i=sm_{j-1}+sa_jx_j$ and so $sa_jx_j\in M_{j-1}$. It follows that $sa_j\in\fq_j$ and so $a_j\in\fq_j$. Therefore $a_jx_j\in M_{j-1}$ and so  $rx_i\in M_{j-1}$. By the inductive hypothesis, we have $(0:_Rx_i+M_{j-1})=\fp_i$ and hence $r\in\fp_i$. Since $r$ is an arbitrary element of $(0:_Rx_i+M_j)$, we obtain $(0:_Rx_i+M_j)\subseteq\fp_i$. The revers inclusion is clear. Therefore  $(0:_Rx_i+M_j)=\fp_i$, as required. This completes the proof of the lemma.
   \end{proof}
 Now, we continue the proof of the theorem.  We set $M_\infty:=\bigcup_{i=1}^\infty M_i$ (note that $M_1\subset M_2\subset M_3\subset \dots$). We show that $(0:_Rx_i+M_\infty)=\fp_i$ for all $i\in\N$. Assume $r\in (0:_Rx_i+M_\infty)$. Therefore $rx_i\in M_j$ for some sufficiently large positive integer $j$. We can assume that $i\leq j$ and so $r\in (0:_Rx_i+M_j)=\fp_i$ by above lemma.  This shows that $(0:_Rx_i+M_\infty)\subseteq\fp_i$. The reverse inclusion is clear and so, in view of the above lemma, $\ass_R(M/M_\infty)$ contains the distinct prime ideals $\fp_1, \fp_2, \dots$, which is impossible because $M$ is weakly Laskerian. Therefore $\supp_R(M/M_i)=\ass_R(M/M_i)$ for some $i$, as required. This proves  the implication (i)$\Rightarrow$(ii) and so the proof of the theorem is completed.
\end{proof}

\end{document}
